\documentclass[nl]{../../elegantbook}
% ========== Omslag en basisinformatie ==========
\title{ElegantBook testdocument}
\subtitle{Kernfuncties van de sjabloon}
\author{Testingenieur}
\institute{Sjabloon Testgroep}
\date{\today}
\version{v1.0}
\extrainfo{Dit document is uitsluitend bedoeld voor sjabloonfunctietesten en heeft geen academische waarde}
\bioinfo{Testbatch}{TEST-FULL-20260506}
\logo{../../figure/logo-blue.png}
\cover{../../figure/cover.jpg}

% ========== Inhoudsopgave en basisinstellingen ==========
\setcounter{tocdepth}{3} % Inhoudsopgave tot niveau 3 weergeven

% ========== Testreferenties ==========
\addbibresource[location=local]{../en/reference.bib}

\begin{document}

\maketitle
\frontmatter
\tableofcontents

% ========== Voorwoord (met sectieonderverdeling) ==========
\pagenumbering{Roman}
\setcounter{page}{0}
\chapter*{Testvoorwoord}
\markboth{Testvoorwoord}{}
\section*{Documentbeschrijving}
\markright{Documentbeschrijving}
Dit document is een volledige testcase voor ElegantBook en omvat alle basis- en geavanceerde functies van de sjabloon. Alle inhoud is fictieve testgegevens zonder praktische betekenis.
\newpage
\section*{Testbereik}
\markright{Testbereik}
Deze test omvat: omslag, inhoudsopgave, lettertypen, meerniveaulijsten, complexe formules, alle wiskundige omgevingen, kruisverwijzingen, afbeeldingen en tabellen, kantnotities, hoofdstuksamenvattingen, oefeningen, referenties en bijlagen.

\mainmatter

% ========== Hoofdstuk 1: Basisopmaak en meerniveaulijsten ==========
\chapter{Basisopmaak en lijsttest}
\begin{introduction}[Testonderdelen van dit hoofdstuk]
\item Chinese lettertypen en tekstopmaak
\item 3-niveau ongenummerde lijstnesting
\item 3-niveau genummerde lijstnesting
\item Gewone alinea's en benadrukte tekst
\end{introduction}

\section{Lettertype en teksttest}
Normale hoofdteksttest. \textbf{Vet}, \textit{Cursief}, \underline{Onderstreept}, \textcolor{red}{Rode teksttest}.

\section{3-niveau ongenummerde lijsttest}
\begin{itemize}
\item Niveau 1 testitem A
\item Niveau 1 testitem B
    \begin{itemize}
    \item Niveau 2 testitem B1
    \item Niveau 2 testitem B2
        \begin{itemize}
        \item Niveau 3 testitem B2-1
        \item Niveau 3 testitem B2-2
            \begin{itemize}
            \item Niveau 4 testitem B2-1a
            \item Niveau 4 testitem B2-2b
            \end{itemize}
        \end{itemize}
    \end{itemize}
\end{itemize}

\section{3-niveau genummerde lijsttest}
\begin{enumerate}
\item Niveau 1 genummerde test 1
\item Niveau 1 genummerde test 2
    \begin{enumerate}
    \item Niveau 2 genummerde test 2-1
    \item Niveau 2 genummerde test 2-2
        \begin{enumerate}
        \item Niveau 3 genummerde test 2-2-1
        \item Niveau 3 genummerde test 2-2-2
            \begin{enumerate}
            \item Niveau 4 genummerde test 2-2-2-1
            \item Niveau 4 genummerde test 2-2-2-2
            \end{enumerate}
        \end{enumerate}
    \end{enumerate}
\end{enumerate}

% ========== Hoofdstuk 2: Test complexe wiskundige formules ==========
\chapter{Test complexe wiskundige formules}
\begin{introduction}
\item Enkelregelige complexe formule
\item Meerregelige formules (align)
\item Matrices / Determinanten
\item Meervoudige integralen / Sommaties / Limieten / Breuken
\item Formulelabels en kruisverwijzingen
\end{introduction}

\section{Enkelregelige complexe formule}
Inline complexe formules: $\sum_{i=1}^n \frac{x_i^2}{\sqrt{y_i+1}} \geq \bar{x}^2$, $\lim_{n\to\infty} \left(1+\frac{1}{n}\right)^n = e$.

\section{Meerregelige formules en matrices}
\begin{align}
a^2 + b^2 &= c^2 \label{eq:gougu} \\
\int_{0}^{1}\int_{0}^{1} f(x,y) dxdy &= 1 \label{eq:double} \\
\begin{vmatrix}
a & b \\
c & d
\end{vmatrix} &= ad-bc \label{eq:det}
\end{align}

\section{Grote operatorformules}
\begin{equation}\label{eq:series}
\sum_{k=0}^{\infty} \frac{x^k}{k!} = e^x,\quad \oint_{C} Pdx+Qdy = \iint_{D}\left(\frac{\partial Q}{\partial x}-\frac{\partial P}{\partial y}\right)dxdy
\end{equation}

Formuleverwijzing: Stelling van Pythagoras \eqref{eq:gougu}, dubbele integraal \eqref{eq:double}, determinant \eqref{eq:det}, reeks \eqref{eq:series}.

% ========== Hoofdstuk 3: Volledige test wiskundige omgevingen ==========
\chapter{Volledige test wiskundige omgevingen}
\begin{introduction}
\item Definitie / Stelling / Lemma / Gevolgtrekking / Axioma / Postulaat
\item Propositie / Voorbeeld / Probleem / Oefening
\item Opmerking / Aanmerking / Conclusie / Aanname / Eigenschap
\item Oplossing / Bewijsomgeving
\end{introduction}

\section{Omgevingstest}

% 1. Definitie
\ifdefined\simpleTest
\begin{definition}[Testdefinitie]\label{def:test}
\else
\begin{definition}[Testdefinitie]{test}
\fi
Dit is de testinhoud van de definitieomgeving, zonder echte wiskundige betekenis, alleen voor stijlverificatie.
\end{definition}

% 2. Stelling
\ifdefined\simpleTest
\begin{theorem}[Teststelling]\label{thm:test}
\else
\begin{theorem}[Teststelling]{test}
\fi
Als $x>0$, dan $x+1>1$, overeenkomend met formule \eqref{eq:gougu}.
\end{theorem}

% 3. Lemma
\ifdefined\simpleTest
\begin{lemma}[Testlemma]\label{lem:test}
\else
\begin{lemma}[Testlemma]{test}
\fi
Voor elk reëel getal $x$ geldt $x^2\geq0$.
\end{lemma}

% 4. Gevolgtrekking
\ifdefined\simpleTest
\begin{corollary}[Testgevolgtrekking]\label{cor:test}
\else
\begin{corollary}[Testgevolgtrekking]{test}
\fi
Uit Lemma \ref{lem:test} volgt: $(x-1)^2\geq0$.
\end{corollary}

% 5. Axioma
\ifdefined\simpleTest
\begin{axiom}[Testaxioma]\label{axi:test}
\else
\begin{axiom}[Testaxioma]{test}
\fi
Testaxioma, zonder praktische betekenis.
\end{axiom}

% 6. Postulaat
\ifdefined\simpleTest
\begin{postulate}[Testpostulaat]\label{pos:test}
\else
\begin{postulate}[Testpostulaat]{test}
\fi
Testpostulaat, alleen voor omgevingsverificatie.
\end{postulate}

% 7. Propositie
\ifdefined\simpleTest
\begin{proposition}[Testpropositie]\label{pro:test}
\else
\begin{proposition}[Testpropositie]{test}
\fi
Testpropositie-inhoud, gekoppeld aan Definitie \ref{def:test}.
\end{proposition}

% 8. Voorbeeld
\begin{example}\label{exa:test}
Voorbeeldomgevingstest, automatische nummering.
\end{example}

% 9. Probleem
\begin{problem}\label{probl:test}
Probleemomgevingstest, voor het stellen van onopgeloste vragen.
\end{problem}

% 10. Oefening
\begin{exercise}\label{exer:test}
Bereken: $1+2+3+\dots+10$.
\end{exercise}

% 11. Opmerking
\begin{note}
Opmerkingsomgeving: zonder nummering, voor belangrijke herinneringen.
\end{note}

% 12. Aanmerking
\begin{remark}
Aanmerkingsomgeving: voor aanvullende uitleg.
\end{remark}

% 13. Conclusie
\begin{conclusion}
Conclusieomgeving: samenvattende testinhoud.
\end{conclusion}

% 14. Aanname
\begin{assumption}
Aannameomgeving: instelling van testvoorwaarden.
\end{assumption}

% 15. Eigenschap
\begin{property}
Eigenschapsomgeving: kenmerken van het testobject.
\end{property}

% 16. Oplossing
\begin{solution}
Oplossing van Oefening \ref{exer:test}: De som is $55$ (testantwoord).
\end{solution}

% 17. Bewijs
\begin{proof}
Bewijs van Stelling \ref{thm:test}: Omdat $x>0$, tel aan beide zijden 1 op; hiermee is de stelling bewezen.
\end{proof}

% ========== Hoofdstuk 4: Kruisverwijzingen · Afbeeldingen · Kantnotities ==========
\chapter{Kruisverwijzingen en uitgebreide functies}
\begin{introduction}
\item Alle soorten kruisverwijzingen
\item Eenvoudige afbeelding- en tabeltest
\item Kantnotitiefunctietest
\item Hoofdstuknavigatieverificatie
\end{introduction}

\section{Samenvatting alle kruisverwijzingen}
Definitie \ref{def:test}, Stelling \ref{thm:test}, Lemma \ref{lem:test}, Gevolgtrekking \ref{cor:test}, Axioma \ref{axi:test}, Postulaat \ref{pos:test}, Propositie \ref{pro:test}, Voorbeeld \ref{exa:test}, Probleem \ref{probl:test}, Oefening \ref{exer:test}.

\section{Afbeelding en tabeltest}
\begin{figure}[h]
    \centering
    \includegraphics[width=0.5\textwidth]{../../image/scatter.jpg}
    \caption{Testafbeelding}\label{fig:test}
\end{figure}
\begin{table}[h]
    \centering
    \begin{tabular}{c*{6}{c}}
    \toprule
        $n$ & 1 & 2 & 3 & 4 & 5 & 6 \\
    \midrule
        $a_n$ & 1 & 1 & 2 & 3 & 5 & 8 \\ 
    \bottomrule
    \end{tabular}
    \caption{Testtabel}\label{tab:test}
\end{table}
Afbeelding/tabelverwijzing: \figref{fig:test} is de testafbeelding, \tabref{tab:test} is de testtabel.

% ========== Hoofdstukafsluitende oefeningen ==========
\begin{problemset}[Testhoofdstukoefeningen]
\item Controleer de weergavestijl van Definitie \ref{def:test}
\item Leid formule \eqref{eq:double} af
\item Controleer of Afbeelding \ref{fig:test} correct wordt weergegeven
\end{problemset}

% ========== Referenties ==========
\nocite{*}
\printbibliography[heading=bibintoc]

% ========== Bijlage ==========
\appendix
\chapter{Testbijlage}
De bijlage-inhoud dient voor verificatie van bijlagehoofdstukken, kopteksten, paginanummering en opmaakstijl, zonder echte informatie.
\section{Testbijlage 1}
De bijlage-inhoud dient voor verificatie van bijlagehoofdstukken, kopteksten, paginanummering en opmaakstijl, zonder echte informatie.
\newpage
\section{Testbijlage 2}
De bijlage-inhoud dient voor verificatie van bijlagehoofdstukken, kopteksten, paginanummering en opmaakstijl, zonder echte informatie.
\end{document}